% controller.tex -- draft of the controller as flown (4 Oct 2026), input from Methodology.tex.
% Notation follows Source/src/model.tex: q_i points from the payload attachment to UAV i and
% T_i = lambda_i q_i. The code (mpc_planner, after Sun et al.) uses s_i = -q_i and t_i = lambda_i.
% Restructured 9 Oct (structure audit): the model moved to model.tex; planner, tracker and steady-error
% rejection are subsections of their own; the result numbers (the box's effect, K's measured values,
% the damping ratio) moved to Section 4.2; the thrust identification procedure moved here from
% hardware.tex; the cost weights moved to Appendix B. Notation: the horizon has N_h nodes over T_h
% (were N and T, which also named the team size and the tension).
% Its references (francis1976internal, pannocchia2003disturbance, pannocchia2007combined,
% maeder2009linear, verschueren2022acados) are in references.bib.

\subsection{Load--cable planner}
\label{sec:controller}

The controller is a two-level cascade of model predictive controllers (Figure~\ref{fig:cascade}).
A centralised planner, based on the load--cable formulation of \cite{tudelft}, plans the motion of
the payload and of every rod over a horizon. Each UAV runs its own tracking MPC, which follows the
planner's reference for that UAV and accounts for the planned rod force. The tracker's output
(collective throttle and three body-rate commands) is sent over the radio link to the flight
controller, whose inner rate loop is not modified. Both levels are solved with acados
\cite{verschueren2022acados} using a real-time iteration (RTI) scheme \cite{diehl2002rti}: one
Gauss--Newton SQP step per tick, warm-started from the previous solution, with the QPs solved by
HPIPM \cite{frison2020hpipm} (partial condensing in the planner, full condensing in the tracker).
Both horizons have $N_h = 20$ nodes over $T_h = 2$\,s.

\begin{figure}[h]
\centering
\begin{tikzpicture}[node distance=1.0cm and 1.3cm, >=Stealth,
    blk/.style={draw, thick, rounded corners=2pt, align=center, minimum height=1.0cm, font=\small},
    lab/.style={font=\scriptsize, align=center}]
  \node[blk, minimum width=3.2cm] (plan) {Load--cable planner\\ \scriptsize 10\,Hz, $N_h=20$, $T_h=2$\,s};
  \node[blk, right=2.0cm of plan, minimum width=2.8cm] (trk) {Tracking MPC $i$\\ \scriptsize 50\,Hz, one per UAV};
  \node[blk, right=1.3cm of trk, minimum width=2.2cm] (fc) {Flight controller\\ \scriptsize rate loop};
  \node[blk, below=1.1cm of trk, minimum width=2.8cm] (sys) {UAVs, rods and payload};
  \node[blk, below=1.1cm of plan, minimum width=3.2cm] (dist) {Payload state, $\mathbf{d}$\\ \scriptsize integral, estimate};
  \draw[->, thick] (plan) -- node[lab, above] {$\mathbf{p}_i, \mathbf{v}_i, \mathbf{a}_i,$\\ $\mathbf{a}_{c,i}, u_{T,i}^{\mathrm{ff}}$} (trk);
  \draw[->, thick] (trk) -- node[lab, above] {$u_T, \boldsymbol{\omega}_{\mathrm{cmd}}$} (fc);
  \draw[->, thick] (fc) |- (sys);
  \draw[->, thick] (sys) -- node[lab, right] {mocap:\\ UAV state} (trk);
  \draw[->, thick] (sys) -- node[lab, below] {mocap:\\ payload pose} (dist);
  \draw[->, thick] (dist) -- node[lab, left] {$\hat{\mathbf{x}}_L, \mathbf{d}$} (plan);
\end{tikzpicture}
\caption[Controller structure]{Controller structure. Motion capture closes both loops; the disturbance term enters the
planner's model.}
\label{fig:cascade}
\end{figure}

The planner works on the model of Section~\ref{sec:model}.

\paragraph{Optimal control problem.} At each planner tick the planner solves
\begin{equation}
\begin{aligned}
    \min_{\mathbf{x}_{0:N_h},\, \mathbf{u}_{0:N_h-1}} \quad & \sum_{k=0}^{N_h-1} \big\| \mathbf{y}(\mathbf{x}_k, \mathbf{u}_k) - \mathbf{y}_k^{\mathrm{ref}} \big\|_{W}^{2}
    + \big\| \mathbf{y}_{N_h}(\mathbf{x}_{N_h}) - \mathbf{y}_{N_h}^{\mathrm{ref}} \big\|_{W_{N_h}}^{2} \\
    \text{s.t.} \quad & \mathbf{x}_{k+1} = \mathbf{F}(\mathbf{x}_k, \mathbf{u}_k; \mathbf{d}), \\
    & C_0 \mathbf{x}_0 = C_0 \hat{\mathbf{x}}(t), \\
    & \lambda_{\min} \le \lambda_{i,k} \le \lambda_{\max}, \quad f_{\min} \le \|\mathbf{f}_{i,k}\| \le f_{\max}, \\
    & |\boldsymbol{\gamma}_{i,k}| \le \gamma_{\max}, \quad |\eta_{i,k}| \le \eta_{\max},
\end{aligned}
\label{eq:plan_ocp}
\end{equation}
with the output
$\mathbf{y} = [\mathbf{p}_L, \mathbf{v}_L, \mathbf{e}_R, \boldsymbol{\omega}_L, \{\mathbf{q}_i\}, \{\lambda_i\}, \{\mathbf{r}_i\}, \mathbf{u}]$,
where $\mathbf{e}_R = 2\,\mathrm{vec}(\mathbf{q}_L^{\mathrm{ref}} \otimes \mathbf{q}_L^{*})$ is the sign-invariant
attitude error. The thrust bounds are softened with slack. $\mathbf{F}$ is the integrated dynamics
over one node of $T_h/N_h = 0.1$\,s.

The initial-condition constraint pins only what motion capture measures directly: the payload pose
and twist, and each rod direction ($C_0$ selects $\mathbf{p}_L, \mathbf{v}_L, \mathbf{q}_L,
\boldsymbol{\omega}_L, \mathbf{q}_i$). The rod rates and their derivatives and the tensions are free
decision variables at node 0, warm-started from the previous solution. Pinning a tension carried over
from the last plan against a measured pose that the UAVs did not quite reach makes node 0 inconsistent.

\paragraph{References.} The payload reference $(\mathbf{p}_L^{\mathrm{ref}}, \mathbf{v}_L^{\mathrm{ref}},
\mathbf{a}_L^{\mathrm{ref}})$ comes from the lift schedule and the trajectory. The rod references are
derived from it by flatness. Holding the payload against gravity, its acceleration and the force
$\mathbf{d}$ needs the rods to support the effective gravity
\begin{equation}
    \mathbf{g}_{\mathrm{eff}} = g\mathbf{e}_3 + \mathbf{a}_{L,xy}^{\mathrm{ref}} - \mathbf{d}/m_L ,
\end{equation}
where $\mathbf{a}_{L,xy}^{\mathrm{ref}}$ is the horizontal part of the reference acceleration (the
trajectories are level),
and the nominal formation (each rod at 45$^\circ$ elevation, tensions $\lambda_i^{0}$ balancing the payload
weight) is rotated so that its vertical axis lies along $\mathbf{g}_{\mathrm{eff}}$, and the tensions are
scaled to $\lambda_i^{\mathrm{ref}} = \lambda_i^{0}\, \|\mathbf{g}_{\mathrm{eff}}\| / g$.

The weights are listed in Table~\ref{tab:weights} (Appendix~\ref{sec:app_params}). The weight on the rod
directions holds the formation's shape, and the weight on the rod rates damps rod swing.

\paragraph{Output to the trackers.} From the solved horizon the planner sends each tracker its
reference position, velocity and acceleration from \eqref{eq:rod_kin}, the planned cable acceleration
$\mathbf{a}_{c,i} = -\lambda_i \mathbf{q}_i / m_q$ at every node, and a throttle feedforward from
$\|\mathbf{f}_i\|$ in \eqref{eq:thrust_req}. The reference is the whole horizon, not a single
set-point, and the tracker interpolates it at its own rate, so a late or skipped planner tick does not
leave it without a reference.

\subsection{UAV tracking MPC and thrust map}
\label{sec:tracker}

Each UAV $i$ runs a tracking MPC on a 17-state model, in which the four stick channels sent to the
flight controller are themselves states and their rates of change are the inputs:
\begin{equation}
    \mathbf{x}_i = [\mathbf{p},\ \mathbf{q},\ \mathbf{v},\ \boldsymbol{\omega},\ \mathbf{c}], \qquad
    \mathbf{c} = [c_\phi, c_\theta, u_T, c_\psi], \qquad \mathbf{u}_i = \dot{\mathbf{c}},
\end{equation}
with
\begin{align}
    \dot{\mathbf{p}} &= \mathbf{v}, &
    \dot{\mathbf{q}} &= \tfrac{1}{2}\, \Omega(\boldsymbol{\omega})\, \mathbf{q}, \\
    \dot{\mathbf{v}} &= R(\mathbf{q})\, k_T u_T \mathbf{e}_3 - g\mathbf{e}_3 - c_z v_z \mathbf{e}_3 + \mathbf{a}_{c,i}, &
    \dot{\boldsymbol{\omega}} &= \frac{1}{\tau_\omega}\big(\boldsymbol{\omega}_{\mathrm{cmd}}(\mathbf{c}) - \boldsymbol{\omega}\big),
    \label{eq:trk_model}
\end{align}
where $\boldsymbol{\omega}_{\mathrm{cmd}}(\mathbf{c})$ is the flight controller's rate curve applied to the
roll, pitch and yaw sticks, and $\tau_\omega$ the time constant of its rate loop. Thrust is modelled
as affine in throttle, identified on the rig: the throttle sent is $u_T + u_0(V)$, with the offset
$u_0(V) = u_0 + \kappa (V_{\mathrm{ref}} - V)$ correcting for the battery voltage $V$ and $k_T$ the gain
above it. The planned cable acceleration $\mathbf{a}_{c,i}$ enters the model as a known, time-varying
parameter along the horizon, so the tracker anticipates the rod's pull instead of reacting to it.
The planned pull is only as good as the planner's model of the rods. On the rig the rods flew at about
60$^\circ$ from horizontal where the plan assumed 45$^\circ$, so each tracker now measures the pull and uses the
plan only for its shape along the horizon: the measured pull is the UAV's acceleration from motion capture
(the least-squares slope of its velocity over 80\,ms, low-pass filtered at 50\,ms) minus the modelled thrust,
gravity and drag, with the thrust passed through the same filter so that a throttle change does not read as
a change of pull. The horizon's planned pull is shifted by the difference between the measured and the
planned pull at its first node, and the result is capped at 12\,m/s$^2$. The planner plans the rods at
60$^\circ$ to match. Section~\ref{sec:res_rig_integral} gives its effect on the rig.

The tracker solves, every 20\,ms,
\begin{equation}
\begin{aligned}
    \min \quad & \sum_{k=0}^{N_h-1} \big\| \mathbf{h}(\mathbf{x}_k, \mathbf{u}_k) - \mathbf{h}_k^{\mathrm{ref}} \big\|_{Q}^{2} + \big\| \mathbf{h}_{N_h}(\mathbf{x}_{N_h}) - \mathbf{h}_{N_h}^{\mathrm{ref}} \big\|_{Q_{N_h}}^{2} \\
    \text{s.t.} \quad & \mathbf{x}_{k+1} = \mathbf{F}_i(\mathbf{x}_k, \mathbf{u}_k; \mathbf{a}_{c,i,k}), \qquad \mathbf{x}_0 = \hat{\mathbf{x}}_i(t), \\
    & 0.05 \le u_{T,k} \le u_{T,\max}, \quad |c_{\phi,k}|, |c_{\theta,k}|, |c_{\psi,k}| \le 1, \quad |\mathbf{u}_k| \le \overline{\mathbf{u}},
\end{aligned}
\label{eq:trk_ocp}
\end{equation}
with $\mathbf{h} = [\mathbf{p}, \mathbf{v}, \boldsymbol{\omega}, \mathbf{c}, \mathbf{u}, \mathbf{e}_R]$, the
attitude reference being the direction of the planned thrust \eqref{eq:thrust_req}. The stick states
are normalised to $[-1, 1]$, and $\overline{\mathbf{u}}$ limits how fast each may change (throttle
0.5\,s$^{-1}$). Only the first command of each solution is sent.

\paragraph{Initial condition.} Node 0 equals the measured state, $\mathbf{x}_0 = \hat{\mathbf{x}}_i(t)$,
which is the standard formulation. The stick states in $\hat{\mathbf{x}}_i$ are the commands actually
sent, not the last solution, so the model starts from where the actuators are. An earlier version
relaxed this to the box $\hat{\mathbf{x}}_i(1-\epsilon) \le \mathbf{x}_0 \le \hat{\mathbf{x}}_i(1+\epsilon)$
with $\epsilon = 0.025$; Section~\ref{sec:res_ic} shows its effect.

\paragraph{Thrust identification.} A proportional map with one gain, identified at hover, did not
hold across loads and battery states. The map was therefore identified on the rig: each UAV hovered
free with hung masses of 0, 98 and 196\,g over battery voltages of 22.7--24.6\,V, and the throttle it
needed was fitted against the mass it carried and the voltage. The fit and $k_T$ are in
Section~\ref{sec:thrust_id}.

\subsection{Steady-error rejection}
\label{sec:disturbance}

Any error that remains constant, from a model mismatch or an external force, is represented by the
constant force $\mathbf{d}$ on the payload in \eqref{eq:plan_trans}. $\mathbf{d}$ is a lumped disturbance:
for an external force (wind, a push, an unmodelled object in the basket) it is that force; for an
internal error (a thrust-map error, the rods' mass) it is the force that would have the same effect on
the payload, which need not exist physically. In the offset-free MPC literature this is an input
disturbance model \cite{pannocchia2003disturbance, maeder2009linear}.

\paragraph{Integral action through the model.} $\mathbf{d}$ is formed from the integral of the payload's
position error $\mathbf{e}_k = \mathbf{p}_L^{\mathrm{ref}} - \mathbf{p}_L$, updated at each planner tick
$k$ (period $T_s = 0.1$\,s):
\begin{equation}
    \mathbf{b}_{k+1} = \mathrm{sat}_{\bar{b}}\!\big(\mathbf{b}_k + k_i T_s\, G_k\, \mathbf{e}_k\big), \qquad
    \mathbf{d}_k = -K \mathbf{b}_k, \qquad K = \mathrm{diag}(K_{xy}, K_{xy}, K_z).
    \label{eq:lumped_int}
\end{equation}
Here $\mathrm{sat}_{\bar{b}}$ bounds each component to $\pm\bar{b}$, and $G_k = \mathrm{diag}(g_{xy}, g_{xy}, g_z)
\in \{0,1\}^3$ is a gate. The gate integrates only in steady, taut flight: lift complete, not
descending or landing, every rod taut, a fresh pose, the payload slow ($<0.05$\,m/s) and the miss
under 0.25\,m. The vertical component also integrates through a slow level orbit. The horizontal
components require a stationary reference, so a value learned in hover is held through a trajectory. A low payload ($e_z > 0$) makes $d_z$ negative: the planner then believes the payload is
heavier and plans more rod tension. Once the payload rests on the floor at landing, $\mathbf{d}$ is
set to zero.

$K$ converts the integral from metres to newtons. It is set to the closed-loop stiffness of the
cascade to a force on the payload, measured by applying a known force and recording the steady
displacement (Section~\ref{sec:res_force}). With this choice the integral's loop gain is
$k_i$, the same as for the equivalent reference-shift form below.

\paragraph{Where $\mathbf{d}$ enters.} $\mathbf{d}_k$ is a parameter of the planner's problem, constant
over the horizon and refreshed at every tick. It enters the payload's translational dynamics
\eqref{eq:plan_trans}, so the plan predicts the payload's motion with the force included and chooses
tensions and UAV positions that counter it. It also enters the reference formation through the
effective gravity $\mathbf{g}_{\mathrm{eff}}$ of the references above, so the formation the cost draws
the solution towards is already tilted and tensioned to balance $\mathbf{d}$, and the cost does not
oppose the force balance.

\paragraph{Zero steady-state error.} Suppose the closed loop converges to an equilibrium at which the
gate is open ($G = I$) and no component of $\mathbf{b}$ is saturated. Then $\mathbf{b}_{k+1} = \mathbf{b}_k$
in \eqref{eq:lumped_int} gives $k_i T_s \mathbf{e} = \mathbf{0}$, so $\mathbf{e} = \mathbf{0}$, whatever
the constant cause. This is the internal model principle \cite{francis1976internal}: a constant
disturbance is rejected because the controller contains a model of a constant, the integrator. The
statement assumes stability; it does not give it. Approximating the cascade's response to $\mathbf{d}$
by a first-order lag with time constant $\tau$, each axis becomes
\begin{equation}
    \tau s^2 + s + k_i = 0, \qquad \zeta = \frac{1}{2\sqrt{k_i \tau}},
\end{equation}
so a faster integral costs damping (Section~\ref{sec:res_force} gives $\zeta$ for the measured $\tau$). Time delays through the planner and the radio link reduce the margin further.

\paragraph{What the integral learns.} Let an unmodelled constant force $\mathbf{F}$ act on the payload,
let $K_c$ be the cascade's actual stiffness to a force on the payload, and treat the cascade as
quasi-static (fast compared with the integral). The plan cancels the force it believes in, $\mathbf{d}$,
so the payload settles displaced by $K_c^{-1}(\mathbf{F} - \mathbf{d})$ and
$\mathbf{e} = -K_c^{-1}(\mathbf{F} - \mathbf{d}) = -K_c^{-1}(\mathbf{F} + K\mathbf{b})$. In continuous time
\eqref{eq:lumped_int} becomes
\begin{equation}
    \dot{\mathbf{b}} = k_i\,\mathbf{e} = -k_i K_c^{-1}\big(K\mathbf{b} + \mathbf{F}\big)
    \quad\Rightarrow\quad \mathbf{b} \to -K^{-1}\mathbf{F}, \qquad \mathbf{d} \to \mathbf{F},
\end{equation}
a first-order convergence at the rate $k_i K K_c^{-1}$ on each axis, with time constant $1/k_i = 2.5$\,s when
$K = K_c$. The believed force converges to the applied one whatever $K$ is: $K$ sets only the speed,
and too large a $K$ costs damping as too large a $k_i$ does. When the force is removed, $\mathbf{b}$
still holds $-K^{-1}\mathbf{F}$, so the payload overshoots by $K_c^{-1}\mathbf{F}$, the displacement the
force caused at onset, and the overshoot decays at the same rate. Section~\ref{sec:res_force} measures both: the
learned force against the applied one, and the overshoot on release.

\paragraph{Reference-shift form.} An earlier version added the height integral to the payload's height
target instead, $z_L^{\mathrm{ref}} \leftarrow z_L^{\mathrm{ref}} + b_z$. Both forms are integral action and
reach the same steady state. In the offset-free MPC literature they correspond to output and input
disturbance models, which differ in transient behaviour, not in the steady state
\cite{pannocchia2007combined}. The model form keeps the target fixed. Its plan predicts a payload
held at the target, which is what happens. The reference-shift form plans a move towards the shifted
target that never comes. The model form also extends to the horizontal axes.

The model form is the default, on the rig as well as in simulation. Where it is not used, the
reference-shift form flies instead, on the height only: wherever the fleet can attach or detach
mid-flight, since the problem is resized with the integral in it, and wherever the integral is
switched off. The rig flights so far used the reference-shift form, as a bounded safety net with
$k_i = 0.4$\,s$^{-1}$.
% STATUS 5 Oct 2026: int_mode auto (model where allowed, else reference): reference with kt_trim on,
% a resizing fleet, z_ki 0, offset_free on or dist_est ring/full (bringup real_mode.planner_options).

\paragraph{Bound.} $\bar{b} = 0.3$\,m bounds the force the integral can supply to $K\bar{b}$ (2.1\,N
vertically), enough for an unmodelled 0.1\,kg object. The bound is a safety measure: a larger error is
a fault, not a disturbance. A warning is raised at 70\,\% of it, so that a fault is reported rather than
silently covered.

\paragraph{Limits.} The integral corrects forces on the payload, not moments: a constant moment, such as
one UAV's thrust-map error, still tilts the payload. A miss larger than the gate is never integrated,
so an error that large stays (Section~\ref{sec:sim_robustness}). The integral is also slow
($1/k_i = 2.5$\,s): it removes constant and slowly varying forces, not impacts or gusts, which only a
force measured on board each UAV could counter quickly.

\paragraph{Alternatives.} Five other ways to remove a steady error were built and flown on the same
cases as the integral (Section~\ref{sec:sim_methods}):
\begin{description}
    \item[Force balance] estimates a constant force on the ring and one shared by the UAVs from the
    measured motion and the delivered thrust, and gives them to the planner's model
    \cite{muske2002disturbance}.
    \item[Offset-free observer] estimates the force from the planner's own prediction error, as in
    offset-free MPC \cite{pannocchia2003disturbance}.
    \item[Measured cable force] each tracker takes its rod's force from its measured acceleration and
    delivered thrust, the ground-side counterpart of the compensation from inertial measurements of
    \cite{tudelft}.
    \item[Thrust-gain trim] a per-UAV estimate of the throttle-to-acceleration gain, from the UAV's own
    measured acceleration and its static share of the rod force.
    \item[Rod mass in the planner] the rods' mass added to the planner's model.
\end{description}

% Under evaluation (sim only): option A, a direct estimate of d from the measured motion and the
% delivered thrust. Per sample, least squares over the tensions, d_L and a force d_D shared by the
% UAVs, from
%   UAV i:    m_q a_i = f_i R_i e3 - m_q g e3 - lambda_i q_i + d_D
%   payload:  m_L a_L = sum_i lambda_i q_i - m_L g e3 + d_L
% filtered by two first-order stages. In the code it is not yet combined with eq:lumped_int.

